\documentclass[wide-notes]{arthur-book}

\title{Arthur Book Minimal Example}
\author{Arthur}

\begin{document}
\maketitle
\tableofcontents

\chapter{A Chapter}
\section{A Section}
\subsection{A Subsection}

This page exercises headings, running headers, margin notes, and common symbols.
\marginkey{margin key}

\begin{arthurnotation}
Use \(\Eps\), \(\Flow\), \(\Ccat\), \(\Pcat\), and \(\Tr\) as core notation.
\end{arthurnotation}

\notationbox

\section{Math}

The equation counter resets by section and prints in the outer margin.

\begin{definition}[Radius]
For a circle in Euclidean space, the radius is the distance from its center to
any point on the boundary.
\end{definition}

\begin{theorem}
If a circle has radius \(r\), then its circumference is \(2\pi r\).
\end{theorem}

\begin{proof}
This is the standard normalization of \(\pi\), written here as a class styling
example.
\end{proof}

\begin{localnotation}
  \notationitem{\abs{x}}{absolute value or size.}
  \notationitem{\set{x}{P(x)}}{the set of \(x\) satisfying \(P(x)\).}
\end{localnotation}

\begin{infobox}[title=Reading note]
Equation summaries are useful when a display introduces several local symbols
that will not be used throughout the whole chapter.
\end{infobox}

\begin{infobox}[margin][title=Margin note]
This is a margin-width info box.
\end{infobox}

\begin{equation}
  C \defeq \underannotate{2\pi r}{circle circumference}
  \label{eq:circumference}
\end{equation}

\equationsummary[margin]{
  \equationterm{\pi}{the ratio between circumference and diameter of a circle.}
  \equationterm{r}{the radius of the circle.}
  \equationterm{C}{the circumference from Equation~\eqref{eq:circumference}.}
}

\begin{break_equation}[eq:long-break]
  F(x)
    &= \alpha_0 + \alpha_1 x + \alpha_2 x^2 + \alpha_3 x^3
      + \alpha_4 x^4 + \alpha_5 x^5
  \eqbreak
      + \alpha_6 x^6 + \alpha_7 x^7 + \alpha_8 x^8 + \alpha_9 x^9 .
\end{break_equation}

\begin{full_equation}[eq:full-width]
  G(x)
    &= \sum_{i=0}^{n} a_i x^i
      + \sum_{j=0}^{m} b_j \exp(-\lambda_j x)
      + \int_{0}^{1} K(x,t)h(t)\,dt
      + \operatorname*{colim}_{c\in\Ccat} H(c)
\end{full_equation}

\begin{cat_math}[eq:pullback][row sep=large,column sep=large]
  P \arrow[r] \arrow[d] & X \arrow[d, "f"] \\
  Y \arrow[r, "g"'] & Z
\end{cat_math}

\clearpage

\section{Figures}

\begin{mainfigure}
\fbox{\rule{0pt}{1.2in}\rule{0.9\textwidth}{0pt}}
\caption{A figure constrained to the main text column.}
\end{mainfigure}

\begin{marginfigure}
\fbox{\rule{0pt}{0.9in}\rule{0.85\marginparwidth}{0pt}}
\caption{A margin-only figure.}
\end{marginfigure}

\begin{fullfigure}
\fbox{\rule{0pt}{1.4in}\rule{0.95\arthurfullwidth}{0pt}}
\caption{A full-width figure spanning the main column and margin column.}
\end{fullfigure}

\end{document}
